\documentclass[10pt]{article}
\usepackage[margin=1in]{geometry}
\usepackage{amsmath,amssymb,amsthm,mathtools}
\usepackage[hidelinks]{hyperref}
\hypersetup{pdftitle={Which unit-gain circulants admit universal perfect state transfer?},pdfauthor={Leo Gao (AI-assisted)}}
\newtheorem{theorem}{Theorem}
\newtheorem{lemma}[theorem]{Lemma}
\newtheorem{corollary}[theorem]{Corollary}
\newcommand{\ZZ}{\mathbb Z}\newcommand{\QQ}{\mathbb Q}
\newcommand{\CC}{\mathbb C}
\title{Which unit-gain circulants admit universal perfect state transfer?}
\author{Leo Gao\\\small AI-assisted research note}
\date{26 November 2026}
\begin{document}
\maketitle
\begin{abstract}
A translation-invariant quantum walk on a ring with unimodular nonzero
hopping amplitudes can transmit a state perfectly between every two sites
only on rings of orders two and three. We parameterize all the triangles
that work by directions of a translated Eisenstein lattice and give the switching-equivalence test.
The key obstruction is that a squared hopping magnitude, after spectral
normalization, must have $p$-adic valuation $-2/\varphi(p^b)$
for every prime power $p^b$ dividing the order; since the magnitude squared is rational,
this valuation must be an integer.
Without the magnitude constraint, we characterize all zero patterns
enabling universal perfect transfer. We also determine which translations
are approximable for generic amplitudes on an arbitrary fixed support.
\end{abstract}

\section{The unit-gain classification}
Write $T e_j=e_{j+1}$ on $\ZZ/n$. A Hermitian circulant Hamiltonian is
$H=\sum_{s=0}^{n-1} h_s T^s$, where $h_{n-s}=\overline{h_s}$ and $h_0$ is real.
A scalar diagonal affects only the global phase, so we subtract it freely. A unit-gain Hamiltonian has $|h_s|=1$ for all its nonzero off-diagonal coefficients.
For a literal property-$\mathbb T$ matrix as in \cite{connelly}, the diagonal must in addition be zero or $\pm1$; simply restore any such diagonal after classification.
Universal perfect state transfer (UPST) means that, for every two vertices
$u,v$, there is a time $t>0$ with $|\langle e_v,e^{-itH}e_u\rangle|=1$.
Here and below we exclude the vacuous one-vertex case.

\begin{theorem}\label{th:main}
A unit-gain Hermitian circulant with UPST has order two or three. In order two the
off-diagonal coefficient is $1$ or $-1$. In order three write
$H=aT+\bar a T^{-1}$, $|a|=1$, and fix $\omega=e^{-2\pi i/3}$. The admissible gains are precisely
\[
 a=\frac{b}{|b|},\qquad b\in (\omega-1)^{-1}+\ZZ[\omega].
\]
Two such loopless triangles are switching equivalent, allowing vertex
permutations, if and only if $a^3=(a')^3$ or $a^3=\overline{(a')^3}$.
In particular, there are infinitely many inequivalent triangles.
\end{theorem}

For context, Cameron et al. \cite{cameron} give the spectral
UPST criterion. Connelly et al. \cite[Sec.~8]{connelly} conjectured
uniqueness of the oriented triangle in the circulant unit-gain setting,
apart from the edge. Their conjecture as literally stated was already
refuted by a three-vertex counterexample in a public formalization report
\cite{truth}; that report explicitly leaves the classification open.
The theorem here is the upper bound on order coupled to an exhaustive
description, not a claim to priority for refuting uniqueness.
The earlier oriented ($\pm i$) theorem \cite{complex} has different
hypotheses. No world-priority claim is made.

\begin{lemma}[Fourier normalization]\label{lem:fourier}
Set $\zeta=e^{2\pi i/n}$ and index eigenvectors so that $T$ has eigenvalue $\zeta^k$.
UPST is equivalent to the existence of $\alpha\in\mathbb R$, $\beta>0$ and
$c_0,\ldots,c_{n-1}\in\ZZ$ for which the eigenvalues in this order are
\[
 \lambda_k=\alpha+\beta(k+n c_k).
\]
For $s\ne0$, setting $z=\zeta^{-s}$, the normalized coefficient is
\[
 \frac{h_s}{\beta}=B(z):=\frac1{z-1}+C(z),\qquad C(X)=\sum_{k=0}^{n-1} c_kX^k.
\]
\end{lemma}
\begin{proof}
Since the propagator commutes with $T$, perfect transfer from $0$ to $-1$
forces $e^{-itH}=\gamma T^{-1}$ for one phase $\gamma$; conversely the powers
of such a propagator perform every translation.
Comparing eigenphases yields the first formula with
$\beta=2\pi/(nt)$. Conversely $t=2\pi/(n\beta)$ works.
Fourier inversion gives the second formula because $\sum z^k=0$ and
$\sum k z^k=n/(z-1)$ at nontrivial $n$th roots of unity (differentiate
the finite geometric sum). No simplicity assumption is needed.
\end{proof}

\begin{lemma}[Prime-power norm obstruction]\label{lem:norm}
Suppose $z$ has order $d=p^r>1$ and $C\in\ZZ[X]$.
For the field norm $N$ from $K=\QQ(z)$, one has $B(z)\ne0$ and
$v_p(N(B(z)))=-1$, with $v_p(p)=1$ on $\QQ$.
If all nonzero $B(w)$ for $w^n=1$, $w\ne1$, have a common absolute value, their common squared magnitude $D$ is rational and
\[
  \varphi(p^r)v_p(D)=-2
       \quad\text{for every }p^r\mid n,\quad r\ge1.
\]
\end{lemma}
\begin{proof}
Let $F(X)=1+(X-1)C(X)$, so $B(z)=F(z)/(z-1)$.
The norm $N(F(z))$ is the resultant of $\Phi_{p^r}$ and $F$.
Modulo $p$, the first polynomial equals $(X-1)^{\varphi(p^r)}$, hence
the resultant is $F(1)^{\varphi(p^r)}=1\pmod p$.
Meanwhile $N(z-1)=\pm\Phi_{p^r}(1)=\pm p$.
This proves the first assertion.
All conjugates of $B(z)\overline{B(z)}$ are $|B(z^u)|^2$ with $u$ coprime to $p$:
complex conjugation commutes with all the cyclotomic automorphisms.
These quantities equal $D$ because none of these conjugates can be zero.
Thus $D$ is fixed by $\mathrm{Gal}(K/\QQ)$, so rational. Taking norms
in $D=B(z)\overline{B(z)}$ gives $D^{\varphi(p^r)}=N(B(z))^2$.
\end{proof}

\begin{proof}[Proof of Theorem~\ref{th:main}]
For a unit-gain Hamiltonian, the common $D$ in the lemma is $1/\beta^2$.
Its valuations are integers. Thus $\varphi(p^r)\mid2$ for every prime
power dividing $n$, leaving only $2,3,4$ as possible prime powers and hence $n\mid 12$.
If $4\mid n$, applying the lemma to $2$ and $4$ gives $v_2(D)=-2$ and
$-1$ simultaneously.
If $6\mid n$, order $2$ gives $D=B(-1)^2$ with $B(-1)\in\QQ$, so $v_3(D)$ is even; order $3$ requires $v_3(D)=-1$.
This eliminates every $n\ge4$.

The two-vertex edge transfers at $\pi/2$.
In order three the set of all $C(\omega)$ in Lemma~\ref{lem:fourier}
is $\ZZ[\omega]$. Normalizing any $b$ in the indicated translate to unit
length makes the off-diagonal coefficients $a,\bar a$. The inverse Fourier construction gives a real diagonal removable without affecting transition magnitudes; it also gives the explicit time $t=2\pi|b|/3$.
The norm lemma ensures $b\ne0$. This proves both directions.
On a triangle the product of edge gains around the oriented cycle is the complete invariant under diagonal unitary switching: match gains on a spanning path using successive vertex phases and the last edge agrees exactly when the products do. Reversing labels conjugates the product. The allowed directions of a translate of a full planar lattice are dense on the unit circle (approximate any ray at distance $R$ by a lattice point within a bounded covering radius and take $R\to\infty$). Thus there are infinitely many admissible directions even after the finite equivalences.
\end{proof}

\begin{corollary}
The same order restriction holds for every finite group $G$ acting regularly
by symmetries of a unit-gain Hamiltonian.
\end{corollary}
\begin{proof}
Choose group coordinates so $H$ commutes with right translations. Transfer
from the identity to $g$ then implies $e^{-itH}=\gamma L_g$, a phase times left multiplication by $g$. In the projective unitary group let $S$ consist of times $t\in\mathbb R$ for which the propagator lies in the finite subgroup of left translations.
This preimage is a closed subgroup of $\mathbb R$, and it is proper:
a continuous image of $\mathbb R$ in a finite group is constant, whereas UPST is nontrivial.
UPST maps $S$ onto the left translations, so $S$ is nonzero. It is therefore of the form $t_0\ZZ$.
Its image is cyclic, forcing $G$ cyclic. Apply Theorem~\ref{th:main}.
\end{proof}

\section{Inverse support problem for exact transfer}
\begin{theorem}\label{th:support}
Given an order $n$, an off-diagonal zero pattern for a Hermitian circulant
with unrestricted complex weights is realizable with UPST if and only if
its zeros are a union of entire additive-order classes, none of prime-power order.
All allowable omissions can be prescribed simultaneously. In particular,
the smallest possible degree is $\sum_{p^r\parallel n}(p^r-1)$, where $p^r$ is the exact power of $p$ dividing $n$.
\end{theorem}
\begin{proof}
By Lemma~\ref{lem:fourier}, zeros away from the diagonal are exactly the
zeros of $F(X)=1+(X-1)C(X)$ among the nontrivial $n$th roots. Since $F$ is integral,
its zeros include all primitive roots of any omitted order.
A prime-power class cannot be omitted by Lemma~\ref{lem:norm}.
Conversely let $J$ be any set of desired omitted orders, all divisors of
$n$ with at least two distinct prime factors. Take
$P(X)=\prod_{d\in J}\Phi_d(X)$. Each $\Phi_d(1)=1$, so $C(X)=(P(X)-1)/(X-1)$ lies in $\ZZ[X]$ and has degree below $n$ (including the zero polynomial if $J=\varnothing$).
The elementary cyclotomic value follows, for example, by induction at $X=1$ in $(X^d-1)/(X-1)=\prod_{1<e\mid d}\Phi_e(X)$.
Our $F$ is exactly $P$ and has precisely the prescribed zeros. Set $\lambda_k=k+n c_k$ with zero-padded coefficients of $C$; its inverse Fourier transform is Hermitian because all eigenvalues are real. To get the degree formula, keep exactly the prime-power classes, which contribute $\sum_{j=1}^r\varphi(p^j)=p^r-1$ per prime.
\end{proof}
This extends the Fourier coefficient formula and composite-order constructions of \cite[Secs.~6--7]{connelly}. It completes the support pattern beyond the prime, prime-square, and power-of-two dense-support cases of \cite{connelly}; all prime-power UPST circulants are dense. It does not impose unit gains.

\section{Approximate versus exact transfer on the same support}
Say that a displacement $m$ admits pretty good state transfer (PGST) when
$\sup_{t>0}|\langle e_m,e^{-itH}e_0\rangle|=1$.
Fix a symmetric set of nonzero allowed steps $S=-S\subset\ZZ/n$.
Write $D_S=\{n/\gcd(n,s):s\in S\}$ and
\[
 M=\prod_{d\in D_S}\Phi_d(1)
   =\prod_{p\mid n}p^{\#\{j\ge1:p^j\in D_S\}} .
\]
In particular $M\mid n$. Here the support specification concerns allowed nonzero step coefficients, not an arbitrary pattern for individual edges.

\begin{theorem}\label{th:generic}
For any Hermitian couplings on $S$, PGST can occur only to
$m\in(n/M)\ZZ/n$. Outside a countable union of proper real linear
hyperplanes in the coupling parameters, it occurs to every such $m$.
The same full-measure conclusion holds on the phase torus for any fixed
positive hopping magnitudes and independent phases on the non-involutive
step pairs (fix either nonzero sign at the unique involution if present).
\end{theorem}
\begin{proof}
By Fourier diagonalization and compactness, transfer with fidelity tending
to one is equivalent to approximating the phases of $T^m$ in every
Fourier coordinate, up to a common phase: the modulus of an average of
unit numbers is one only when all its summands agree.
The character form of Kronecker's theorem for a continuous torus flow,
applied to the differences $\lambda_k-\lambda_0$, therefore says that $m$ works precisely when
\[
 n\mid m f'(1)\quad\hbox{for all }f(X)=\sum_{k=0}^{n-1} l_k X^k
 \text{ with }l_k\in\ZZ,\ f(1)=0,\ \sum_k l_k\lambda_k=0 .
\]
Indeed the character indexed by $(l_k)$ kills the flow modulo common phase under just these two zero conditions and it evaluates on the target as $\zeta^{m\sum k l_k}$.
Positive times have the same closure as all real times by recurrence of a compact torus flow.

Relations independent of coupling values are precisely those with
$f(1)=0$ and $f(\zeta^s)=0$ for every $s\in S$. In fact the relation is
$\sum_{s\in S}h_s f(\zeta^s)=0$; each step pair contributes
$2\Re(h_s f(\zeta^s))$, and an involution contributes a real scalar.
It vanishes identically precisely when every coefficient is zero. For any
of the countably many other $f$ it defines a proper real hyperplane.
For fixed magnitudes it instead defines a nonidentically-zero real trigonometric polynomial in the phases: the distinct one-coordinate Fourier modes cannot cancel, and the only possible constant term comes from the lone involution. Such a zero set has measure zero.

Universal relations are exactly the multiples of
$L(X)=(X-1)\prod_{d\in D_S}\Phi_d(X)$ of degree below $n$.
All give $f'(1)$ divisible by $M=L'(1)$. For necessity one may reduce
$L$ modulo $X^n-1$: every $n$th-root evaluation stays unchanged, and its derivative at $1$ changes by a multiple of $n$. Thus the test above implies $n\mid mM$ even in the edge case $L=X^n-1$. Conversely when no additional relations hold, $n\mid mM$ suffices.
\end{proof}

A nearest-neighbor ring with $n$ having two distinct prime factors
cannot have nontrivial PGST at \emph{any} magnetic flux. On $p^r$ sites,
exactly $p$ sites (counting the original site) are generically reachable;
this says nothing by itself about the waiting time. Universal PGST
needs only one step pair from every prime-power order dividing $n$,
in contrast to the entire order classes required for UPST. Earlier
papers give oriented prime cycles \cite{cameron} and real unweighted
cycles \cite{pal} as special (not generic) cases.

\paragraph{Reproducibility and boundaries.}
The companion \texttt{verify.py} checks prescribed supports using the
inverse Fourier polynomial directly; it checks prime-power norms by exact
integer determinants, and checks explicit triangle parameters. None of
the proofs relies on these finite tests. The results do not impose an upper
bound on arbitrary noncirculant unit-gain Hamiltonians with no regular
group symmetry. Existence of PGST is not a useful-time transfer estimate.

\begin{thebibliography}{9}
\bibitem{cameron} S.~Cameron, S.~Fehrenbach, L.~Granger, O.~Hennigh, S.~Shrestha, and C.~Tamon,
\emph{Universal state transfer on graphs}, Linear Algebra Appl.~455 (2014), 115--142.
\href{https://arxiv.org/abs/1310.3885}{arXiv:1310.3885}.
\bibitem{connelly} E.~Connelly, N.~Grammel, M.~Kraut, L.~Serazo, and C.~Tamon,
\emph{Universality in perfect state transfer}, Linear Algebra Appl.~531 (2017), 516--532.
\href{https://arxiv.org/abs/1701.04145}{arXiv:1701.04145};
\href{https://doi.org/10.1016/j.laa.2017.06.015}{journal DOI}.
\bibitem{complex} A.~Acuaviva, A.~Chan, S.~Eldridge, C.~Godsil, M.~How-Chun-Lun, C.~Tamon, E.~Wright, and X.~Zhang, \emph{State transfer in complex quantum walks}, Philos. Trans. Roy. Soc. A 383 (2025), 20240420,
\href{https://arxiv.org/abs/2301.01473}{arXiv:2301.01473}.
\bibitem{pal} H.~Pal and B.~Bhattacharjya, \emph{Pretty good state transfer on circulant graphs},
Electronic J. Combin. 24 (2017), P2.23.
\href{https://arxiv.org/abs/1607.03598}{arXiv:1607.03598}.
\bibitem{truth} trureturing, \emph{A unit-gain circulant on three vertices with universal perfect state transfer}, formalization report with explicit order-three witness, accessed 26 November 2026.
\href{https://www.trureturing.org/research/connelly-2017-property-t-universal-state-transfer/}{Primary project report}.
\end{thebibliography}
\end{document}
